Praca dotyczy projektu i analizy wydajności komunikacji bezprzewodowej w systemach czasu rzeczywistego z wykorzystaniem mikrokontrolera ESP32. W ramach realizacji opracowano autorskie środowisko badawcze ESPy-Nexus w architekturze Hardware-in-the-Loop, łączące rzeczywiste urządzenie wbudowane z komputerowym silnikiem testowym, bazą danych oraz webowym interfejsem analitycznym. Celem badania było empiryczne określenie wpływu protokołu komunikacyjnego, topologii sieci, rozmiaru ładunku i częstotliwości transmisji na niezawodność, terminowość i przepustowość przesyłania danych. W pracy zaprojektowano i zautomatyzowano zestaw eksperymentów obejmujących protokoły SERIAL, UDP, TCP i WebSocket oraz różne konfiguracje sieciowe. Zbierane dane przetwarzano do wskaźników QoS, takich jak PDR, jitter, goodput, czas przerwy w odbiorze, względne opóźnienie oraz dryft zegara. Wyniki pokazały, że nie istnieje jeden dominujący protokół: UDP zapewnia lepszą responsywność i wysoką przepustowość kosztem strat pakietów, natomiast TCP i WebSocket lepiej zachowują integralność strumienia, ale mogą powodować duże opóźnienia przy dużym obciążeniu. Opracowanie stanowi praktyczny fundament do dalszych badań nad doborem protokołów komunikacyjnych w systemach sterowania i telemetrii czasu rzeczywistego.

Słowa kluczowe: Python; ESP32; komunikacja bezprzewodowa; QoS; UDP; TCP; WebSocket; Hardware-in-the-Loop; sieci Wi-Fi; czas rzeczywisty; telemetryka; wydajność sieciowa
